\subsection{Federated Learning Pipeline}
\label{sec:fl_architecture}

The Federated Learning pipeline enables collaborative model training
across multiple nodes without sharing their raw data, 
ensuring data privacy and distributing the computational
load. The pipeline was implemented using the \textit{Flower} framework
\cite{Beuteletal2020} and its \textit{Virtual Client Engine (VCE)},
which enables experimentation with multiple federation setups. In this
mini-project, the simulation setup was configured as follows:

\begin{itemize}
\item \textbf{Clients:} 4 independent nodes.
\item \textbf{Data split:} Randomly distributed (Independent and Identically
  Distributed (IID), 25\% of the dataset per client).
\item \textbf{Training:}
  \begin{itemize}
  \item \textbf{Local Execution:} 1 epoch per client per round.
  \item \textbf{Global Execution:} 50 total communication rounds
    (equivalent to the centralized baseline epochs).
  \item \textbf{Synchronization:} Synchronous updates (nodes wait for
    all others to finish before starting the next round).
  \end{itemize}

 \item \textbf{Model Aggregation Strategy:} \texttt{FedAvg}
\end{itemize}

A single local execution epoch was configured to avoid local model
overfitting, which commonly occurs when a high number of local epochs are
executed within the federation. This ensures that one federation round is
equivalent to one centralized epoch in terms of data passes through
the model. The standard \texttt{FedAvg} aggregation algorithm was used
to perform the global model weights update.